\subsection{Projective dimension, injective dimension, homological dimension}

Let A be a ring, and M an A-module. Recall (A, X, p.~134, def. 1) that the projective dimension of M, denoted \(\mathrm{dp}_{\mathrm{A}}(\mathrm{M})\), is the lower bound (in \(\overline{\mathbf{Z}}\) ) of the set of lengths of projective resolutions of M. We have \(\mathrm{dp}_{\mathrm{A}}(0) = -\infty\) and \(\mathrm{dp}_{\mathrm{A}}(\mathrm{M}) \geqslant 0\) if M is nonzero. For M to be projective, it is necessary and sufficient that one has \(\mathrm{dp}_{A}(\mathrm{M}) \leqslant 0\) .

Example. Let J be an ideal of A generated by an A-regular sequence \(\mathbf{x} = (x_{1}, \ldots, x_{r})\) . We have dp \(_{A}\) (A/J) ≤ r. Indeed, this is clear if A is zero; otherwise, the Koszul complex \(\mathbf{K}_{\bullet}(\mathbf{x}, \mathrm{A})\) is a free resolution of A/J of length r (loc. cit., p.~159, remark 3). Moreover, for every A-module N, the A-modules Ext \(_{A}^{r}\) (A/J, N) and N/JN are isomorphic (loc. cit.); consequently, for one to have dp \(_{A}\) (A/J) = r, it is necessary and sufficient that J be distinct from A (loc. cit., p.~134, prop. 1).

One defines similarly the injective dimension of M, denoted \(\mathrm{di}_{\mathrm{A}}(\mathrm{M})\) , as the lower bound of the set of lengths of injective resolutions of M. We have \(\mathrm{di}_{\mathrm{A}}(0) = -\infty\) , and \(\mathrm{di}_{\mathrm{A}}(\mathrm{M}) \geqslant 0\) if \(M \neq 0\) . For M to be injective, it is necessary and sufficient that one has \(\mathrm{di}_{\mathrm{A}}(\mathrm{M}) \leqslant 0\) .

PROPOSITION 1.- Let A be a ring, M an A-module, and n an integer \(\geqslant 0\). The following conditions are equivalent:

\begin{enumerate}
\def\labelenumi{(\roman{enumi})}
\item
  we have \(\mathrm{di}_{\mathrm{A}}(\mathrm{M}) \leqslant n\) (in other words, M possesses an injective resolution of length \(\leqslant n\) );
\item
  for every A-module N and every integer \(i > n\) , one has \(\mathrm{Ext}_A^i (\mathrm{N},\mathrm{M}) = 0\) ;
\item
  for every ideal \(\mathfrak{a}\) of A, we have \(\operatorname{Ext}_{\mathrm{A}}^{n+1}(\mathrm{A}/\mathfrak{a},\mathrm{M})=0\) ;
\item
  for every exact sequence of A-modules
\end{enumerate}

\[
0 \rightarrow \mathrm{M} \longrightarrow \mathrm{I} ^ {0} \longrightarrow \mathrm{I} ^ {1} \longrightarrow \dots \longrightarrow \mathrm{I} ^ {n - 1} \longrightarrow \mathrm{Q} \rightarrow 0,
\]

where the \(I^{i}\) are injective, the A-module Q is injective.

(i)⇒(ii): this follows from A, X, p.~100, th. 1.

(ii)⇒(iii): This is clear.

\((\mathrm{iii}) \Rightarrow (\mathrm{iv}) : \text{in the situation of (iv), for every A-module N there is an isomorphism of Ext}_A^1(\mathrm{N}, \mathrm{Q}) \text{onto Ext}_A^{n+1}(\mathrm{N}, \mathrm{M}) (\mathrm{A}, \mathrm{X}, \text{p. 128, cor. 4}) ; \text{under the hypothesis (iii), the A-module Ext}_A^1(A / \mathfrak{a}, \mathrm{Q}) \text{ is zero for every ideal } \mathfrak{a} \text{ of A and Q is injective (A, X, p. 93, prop. 11).}\)

(iv)⇒(i): Consider the exact sequence (A, X, p.~52)

\[
0 \longrightarrow \mathrm{M} \longrightarrow \mathrm{I} ^ {0} (\mathrm{M}) \longrightarrow \dots \longrightarrow \mathrm{I} ^ {n - 1} (\mathrm{M}) \longrightarrow \mathrm{K} ^ {n - 1} (\mathrm{M}) \longrightarrow 0;
\]

if condition (iv) is satisfied, the A-module \(K^{n-1}(M)\) is injective, whence (i).

Recall (A, X, p.~138, def. 2) that the homological dimension of the ring A, denoted dh(A), is the supremum in \(\overline{\mathbf{Z}}\) of the set of integers \(n\) for which there exist two A-modules M and N such that Ext \(_{A}^{n}\) (M,N) is nonzero. It is therefore also the supremum of the set of projective (or injective) dimensions of all A-modules; when A is Noetherian, one may restrict to finitely generated A-modules (A, X, p.~139, cor.).

